\documentclass[11pt]{article}

\usepackage[margin=1in]{geometry}
\usepackage[T1]{fontenc}
\usepackage[utf8]{inputenc}
\usepackage{lmodern}
\usepackage{microtype}
\usepackage{amsmath,amssymb,amsthm}
\usepackage{mathtools}
\usepackage{stmaryrd}
\usepackage{xcolor}
\usepackage{hyperref}
\usepackage{booktabs}
\usepackage{enumitem}
\usepackage{listings}

\hypersetup{
  colorlinks=true,
  linkcolor=blue!50!black,
  citecolor=blue!50!black,
  urlcolor=blue!50!black
}

\lstdefinelanguage{Lean4}{
  morekeywords={
    import,namespace,end,open,section,noncomputable,abbrev,structure,inductive,
    def,theorem,lemma,where,match,with,if,then,else,by,fun,let,in,return,do,
    simp,calc,have,show,exact,rcases,cases,constructor,intro,apply,
    deriving,instance,class,extends,private,protected,mutual,partial,
    Type,Prop,Sort,Nat,Bool,true,false,sorry,axiom
  },
  sensitive=true,
  morecomment=[l]{--},
  morecomment=[s]{/-}{-/},
  morestring=[b]"
}

\lstset{
  language=Lean4,
  basicstyle=\ttfamily\small,
  keywordstyle=\bfseries\color{blue!50!black},
  commentstyle=\itshape\color{green!40!black},
  stringstyle=\color{red!50!black},
  showstringspaces=false,
  breaklines=true,
  frame=single,
  framerule=0.3pt,
  xleftmargin=0.5em,
  xrightmargin=0.5em,
  keepspaces=true,
  columns=fullflexible
}

% Theorem environments
\newtheorem{theorem}{Theorem}[section]
\newtheorem{lemma}[theorem]{Lemma}
\newtheorem{proposition}[theorem]{Proposition}
\newtheorem{corollary}[theorem]{Corollary}
\newtheorem{conjecture}[theorem]{Conjecture}
\theoremstyle{definition}
\newtheorem{definition}[theorem]{Definition}
\newtheorem{example}[theorem]{Example}
\theoremstyle{remark}
\newtheorem{remark}[theorem]{Remark}

% Quantale operators (stmaryrd provides \bigsqcup; define \bigsqcap if missing)
\providecommand{\bigsqcap}{\mathop{\textstyle\sqcap}\limits}

\title{Process Calculi Formalized in Lean~4:\\
\large Rho, Pi, MQ, and Modal Mu-Calculus \\
with Cross-Calculus Bridges}
\author{Zar \and Oru\v{z}i\footnote{%
``Oru\v{z}i'' denotes AI-assisted collaboration
using Claude Code (Anthropic), ChatGPT, and Codex (OpenAI).}}
\date{March 2026}

\begin{document}
\maketitle

\begin{abstract}
We present a comprehensive Lean~4 formalization of five process-algebraic
systems---the $\rho$-calculus, $\pi$-calculus, MQ-calculus,
MeTTa-calculus, and modal $\mu$-calculus---together with three
cross-calculus bridges: a $\pi$-to-$\rho$ encoding with forward
simulation and backward admin reflection, an MQ-to-MORK
interoperability bridge, and a MeTTa-to-$\rho$ shared-fragment
translation with open-map self-bisimulation witness.
The formalization totals over 17{,}500 lines of Lean~4.28 across 55
files with zero \texttt{sorry} and zero custom axioms.
Shared infrastructure---typeclasses for parallel composition,
restriction, structural congruence, and reflexive-transitive
closure---eliminates approximately 40\% of cross-calculus duplication.
Quantale-valued semantics for $\mu$-calculus enables instantiation to
probabilistic, fuzzy, tropical, and PLN-evidence domains.
\end{abstract}


%% ====================================================================
\section{Introduction}
\label{sec:intro}
%% ====================================================================

Process calculi provide compositional models of concurrent, communicating
systems.  In this paper we report on the formal verification, in the
Lean~4 proof assistant~\cite{Moura2021}, of five such calculi and the
bridges that connect them.

\begin{enumerate}[leftmargin=1.5em]
  \item The \textbf{$\rho$-calculus}~\cite{MeredithRadestock2005}: a
    reflective higher-order calculus where names \emph{are} quoted
    processes, supporting structural congruence, a COMM reduction rule,
    an executable engine with soundness proof, the Spice
    calculus~\cite{MeredithPresentMoment2026} for $n$-step lookahead, and
    derived replication and restriction.
  \item The \textbf{$\pi$-calculus}~\cite{Milner1999,SangiorgiWalker2001}:
    the standard calculus of mobile processes, with a full encoding into
    the $\rho$-calculus following Lybech's name-server
    approach~\cite{Lybech2022}, including forward simulation, encoding
    morphism structure, weak bisimilarity, and backward admin reflection.
  \item The \textbf{MQ-calculus}~\cite{StayMeredithMQ2026}: a quantum
    process calculus where communication \emph{is} measurement, producing
    binary outcomes governed by Born-rule probabilities.  We formalize
    backend-parametric denotational semantics with a statevector
    instantiation and an interoperability bridge to
    MORK~\cite{MeredithStayOSLF2020}.
  \item The \textbf{MeTTa-calculus}~\cite{MeredithMeTTaCalculus2024}: a
    symmetric reflective higher-order concurrent calculus with
    backchaining, formalizing the core of the MeTTa language with COMM
    and REFL rules, an executable reduction engine, a premise-IR
    adequacy bridge, and a shared-fragment translation into the
    $\rho$-calculus with open-map self-bisimulation witness.
  \item The \textbf{modal $\mu$-calculus}~\cite{Kozen1983,BradfieldStirling2007}:
    propositional modal logic with least and greatest fixed points.  We
    provide both standard Boolean semantics and quantale-valued
    semantics, and prove that the Galois connection in
    $\rho$-calculus \emph{is} diamond-box duality.
\end{enumerate}

Three bridges tie these together: the $\pi$-to-$\rho$ encoding with
simulation theorems, an MQ-MORK binary-fold bridge, and a
MeTTa-to-$\rho$ shared-fragment translation.
All proofs are checked by the Lean~4.28 kernel with zero
\texttt{sorry}~\cite{mathlib2020}.

Table~\ref{tab:summary} gives line counts; total formalization exceeds
17{,}500 lines.

\begin{table}[t]
\centering
\small
\begin{tabular}{llrr}
\toprule
\textbf{Component} & \textbf{Location} & \textbf{Files} & \textbf{Lines} \\
\midrule
$\rho$-calculus  & \texttt{ProcessCalculi/RhoCalculus/} & 11 & 4{,}954 \\
$\pi$-calculus   & \texttt{ProcessCalculi/PiCalculus/}  & 19 & 9{,}134 \\
MQ-calculus      & \texttt{ProcessCalculi/MQCalculus/}  & 10 & 1{,}433 \\
MeTTa-calculus   & \texttt{ProcessCalculi/MeTTaCalculus/} &  9 &   904 \\
Common infra     & \texttt{ProcessCalculi/Common/}      &  4 &   246 \\
Modal $\mu$-calc & \texttt{Logic/ModalMuCalculus.lean}   &  1 &   360 \\
Quantale sem.    & \texttt{Logic/ModalQuantaleSemantics.lean} & 1 & 520 \\
\midrule
\textbf{Grand total} & & \textbf{55} & \textbf{17{,}551} \\
\bottomrule
\end{tabular}
\caption{Formalization summary: 0 sorries, 0 custom axioms throughout.}
\label{tab:summary}
\end{table}


%% ====================================================================
\section{The Rho-Calculus}
\label{sec:rho}
%% ====================================================================

The $\rho$-calculus~\cite{MeredithRadestock2005} is a reflective
higher-order process calculus in which names are quoted processes:
every name $x$ is of the form $\mathsf{@}P$ for some process~$P$.
This ``reflection'' dissolves the traditional name/process distinction
and raises expressiveness questions beyond purely behavioral
formalisms (Section~\ref{sec:mu-bridge}).
The formalization comprises 11~files totaling 4{,}954~lines under
\texttt{ProcessCalculi/RhoCalculus/}.

\subsection{Syntax}

Process syntax uses a locally nameless representation.  The core
inductive type \texttt{Pattern} (in \texttt{Types.lean}, 221~lines)
includes:

\begin{lstlisting}[caption={Rho-calculus process syntax (from \texttt{Types.lean})}]
inductive Pattern : Type where
  | pNil    : Pattern
  | pPar    : Pattern -> Pattern -> Pattern
  | pInput  : Pattern -> Pattern -> Pattern -> Pattern
  | pOutput : Pattern -> Pattern -> Pattern
  | pDrop   : Pattern -> Pattern
  | pQuote  : Pattern -> Pattern
  | pVar    : Nat -> Pattern
  deriving DecidableEq, Repr
\end{lstlisting}

\noindent
The \texttt{pQuote}/\texttt{pDrop} pair implements reflection: quoting a
process produces a name, and dropping a name recovers the process.

\subsection{Structural Congruence}

Structural congruence (\texttt{StructuralCongruence.lean}, 287~lines)
is formalized as an inductive relation including alpha-equivalence,
parallel commutativity and associativity, flattening of nested parallels,
and the property that quoting respects structural congruence:

\begin{theorem}[Quote Respects SC]
If $P \equiv_{\mathrm{SC}} Q$ then $\mathsf{@}P \equiv_{\mathrm{SC}} \mathsf{@}Q$.
\end{theorem}

\subsection{COMM Rule and Reduction}

The COMM rule (\texttt{CommRule.lean}, 298~lines) formalizes the
core communication step: when an output $\bar{x}\langle y \rangle$
meets an input $x(z).P$, the name $y$ is substituted for~$z$ in~$P$.
The full reduction relation (\texttt{Reduction.lean}, 430~lines) closes
COMM under parallel contexts and structural congruence.

\subsection{Executable Engine and Soundness}

An executable reduction engine (\texttt{Engine.lean}, 427~lines)
provides a fuel-bounded decision procedure for single-step reduction.
A separate soundness proof (\texttt{Soundness.lean}, 739~lines) shows
that every engine-computed reduction is derivable in the inductive
relation:

\begin{theorem}[Engine Soundness]
If \texttt{engineStep fuel p = some q} then \texttt{Reduces p q}.
\end{theorem}

\subsection{Spice Calculus: Present Moment and Lookahead}

The Spice calculus~\cite{MeredithPresentMoment2026} extends $\rho$-calculus
COMM with $n$-step lookahead (\texttt{SpiceRule.lean}, 419~lines;
\texttt{PresentMoment.lean}, 923~lines).  Agents receive the set of
states reachable in ${\leq}\, n$ steps before committing to an
interaction:
\[
  \texttt{futureStates}(p, n) = \{q \mid p \leadsto^n q\}
\]
When $n = 0$ this recovers standard $\rho$-calculus COMM.  The ``present
moment'' ($n = 1$) gives agents awareness of their immediate interaction
possibilities.  Key theorems include compositional $n$-step parallel
congruence lifting (\texttt{ReducesN.par\_head},
\texttt{ReducesN.par\_any\_pos}) and
\texttt{futureStates\_par\_head}.

\subsection{Derived Replication and Restriction}

The file \texttt{DerivedRepNu.lean} (669~lines) derives replication
$!P$ and restriction $(\nu x)\, P$ as definable constructs within
the $\rho$-calculus, proving they satisfy the expected reduction and
congruence laws.


%% ====================================================================
\section{The Pi-Calculus and the Pi-to-Rho Encoding}
\label{sec:pi}
%% ====================================================================

The $\pi$-calculus~\cite{Milner1999,SangiorgiWalker2001} is formalized
in 19~files totaling 8{,}967~lines under \texttt{ProcessCalculi/PiCalculus/}.
We follow the Lybech 2022 encoding~\cite{Lybech2022} of $\pi$-calculus
into $\rho$-calculus via a name-server approach.

\subsection{Syntax and Reduction}

The $\pi$-calculus syntax (\texttt{Syntax.lean}, 85~lines) and
reduction relation (\texttt{Reduction.lean}, 405~lines) follow
standard presentations~\cite{SangiorgiWalker2001}.

\subsection{Lybech 2022 Encoding}

The encoding (\texttt{RhoEncoding.lean}, 375~lines) maps $\pi$-calculus
names to $\rho$-calculus names via a name server that mediates
between the flat name space of $\pi$-calculus and the reflective name
space of $\rho$-calculus.

\subsection{Forward Simulation}

The key correctness result is forward simulation
(\texttt{ForwardSimulation.lean}, 1{,}387~lines): if a $\pi$-calculus
process reduces, its encoding reduces (modulo administrative steps) in
$\rho$-calculus.

\begin{theorem}[Forward Simulation]
\label{thm:forward-sim}
If $P \to_\pi P'$ then $\llbracket P \rrbracket \to_\rho^* \llbracket P' \rrbracket$
(up to administrative reductions).
\end{theorem}

\noindent
The theorem name in Lean is \texttt{encoding\_forward\_single\_step\_rf}.

\subsection{Encoding Morphism Structure}

The file \texttt{EncodingMorphism.lean} (894~lines) establishes that
the encoding preserves the algebraic structure of parallel composition,
restriction, and nil:
\[
  \llbracket P \mathbin{|} Q \rrbracket \equiv \llbracket P \rrbracket \mathbin{|} \llbracket Q \rrbracket, \qquad
  \llbracket \mathbf{0} \rrbracket \equiv \mathbf{0}, \qquad
  \llbracket (\nu a)\, P \rrbracket \equiv (\nu \hat{a})\, \llbracket P \rrbracket
\]

\subsection{Weak Bisimilarity}

Weak bisimilarity (\texttt{WeakBisim.lean}, 230~lines;
\texttt{WeakBisimDerived.lean}, 149~lines) formalizes the standard
notion from~\cite{SangiorgiWalker2001}: two processes are weakly
bisimilar when every observable reduction step of one can be matched
(up to silent $\tau$-steps) by the other, and vice versa.
The definition is Prop-valued and coinductive-style: a witness is an
explicit relation $R$ that is symmetric, closed under forward and
backward lifting, and preserves barbed observability restricted to a
finite name set~$N$ (following Lybech's $N$-restricted barbed
bisimilarity~\cite{Lybech2022}).
\texttt{WeakBisimOpenMapBridge.lean} connects this to the open-maps
approach of Joyal, Nielsen, and Winskel~\cite{JoyalNielsenWinskel1996}
by showing that weak $N$-restricted bisimilarity coincides with
path bisimilarity over a \texttt{BisimulationKit} that instantiates the
generalized open-map framework from
\texttt{CategoryTheory/GeneralizedOpenMaps.lean}.

\subsection{Backward Admin Reflection}

The largest single file in the formalization,
\texttt{BackwardAdminReflection.lean} (2{,}745~lines), proves that
administrative reductions introduced by the encoding can be reflected
backward: if $\llbracket P \rrbracket \to_\rho^* Q$ via admin steps,
then $Q$ is structurally congruent to some $\llbracket P' \rrbracket$.
The file now also exposes a fully general backward-simulation endpoint over
encoded step-sources:
\texttt{weak\_backward\_full\_simulation\_encodedSC}, together with encoded-source
specializations
\texttt{weak\_backward\_full\_simulation\_encode} and
\texttt{weak\_backward\_full\_simulation\_encode\_rfStepAt}.
These theorems flatten the reflected trace decomposition into three explicit
branches: reflexive endpoint, admin-progress branch, or reflected $\pi$-step
branch with head/tail trace data.  This completes the backward direction of
the simulation argument: every administrative reduction sequence starting
from an encoded process eventually reaches an encoded image again.


%% ====================================================================
\section{The MQ-Calculus}
\label{sec:mq}
%% ====================================================================

The MQ-calculus~\cite{StayMeredithMQ2026} unifies communication and
quantum measurement: a COMM event produces a binary measurement outcome
governed by Born-rule probabilities.  The formalization
(10~files, 1{,}433~lines) covers syntax, operational reduction,
backend-parametric denotational semantics, and a cross-calculus bridge
to MORK.
This extends the communicating quantum processes
framework of Gay and Nagarajan~\cite{GayNagarajan2005} with explicit
measurement branching in the reduction semantics.

\subsection{Syntax}

The process grammar uses de~Bruijn-style wire indexing and typed gate
specifications:

\begin{lstlisting}[caption={MQ-calculus process syntax (from \texttt{Syntax.lean})}]
inductive GateOp : Type where
  | H | X | Z
  | custom (name : String)
  deriving DecidableEq, Repr

structure GateSpec where
  op : GateOp
  target : Nat
  deriving DecidableEq, Repr

inductive Process : Type where
  | MQNil  : Process
  | MQPar  : Process -> Process -> Process
  | MQNu   : Process -> Process
  | MQGate : GateSpec -> Process -> Process
  | MQOut  : Nat -> Process
  | MQIn   : Nat -> Process -> Process -> Process
  deriving DecidableEq, Repr
\end{lstlisting}

\noindent
The \texttt{MQIn} constructor takes two continuations---one for each
measurement outcome---making measurement branching explicit in the
syntax.

\subsection{Communication as Measurement: The COMM Rule}

The COMM rule (\texttt{CommRule.lean}, 87~lines) produces a binary
measurement outcome.  A central theorem provides a constructive
both-outcomes witness:

\begin{lstlisting}[caption={COMM rule with constructive both-branches (from \texttt{CommRule.lean})}]
inductive Outcome : Type where
  | zero : Outcome
  | one  : Outcome
  deriving DecidableEq

inductive CommReduction :
    Nat -> Process -> Process -> MeasurementBranch -> Prop where
  | outcome_zero (i : Nat) (p q : Process) :
      CommReduction i p q { outcome := .zero, result := p }
  | outcome_one (i : Nat) (p q : Process) :
      CommReduction i p q { outcome := .one, result := q }

theorem comm_both_outcomes (i : Nat) (p q : Process) :
    CommReduction i p q { outcome := .zero, result := p }
    /\ CommReduction i p q { outcome := .one, result := q } :=
  (.outcome_zero i p q, .outcome_one i p q)
\end{lstlisting}

\noindent
Born-style probabilities are defined separately on quantum state
vectors; the branch relation itself remains structural and
proof-friendly.

\subsection{Backend-Parametric Denotational Semantics}

The semantics (\texttt{Backend.lean}, 380~lines;
\texttt{Denotational.lean}, 229~lines) is parameterized by a backend
interface that provides gate action, fresh allocation, branch
probabilities, and post-measurement collapse:

\begin{lstlisting}[caption={Backend interface (from \texttt{Backend.lean})}]
structure MQSemanticsBackend where
  branchProb      : Nat -> Outcome -> QState -> Real
  branchProb_nonneg : forall i out st, 0 <= branchProb i out st
  branchProb_sum_one : forall i st,
    branchProb i .zero st + branchProb i .one st = 1
  applyGate       : GateSpec -> QState -> QState
  allocFresh      : QState -> QState
  collapse        : Nat -> Outcome -> QState -> QState
\end{lstlisting}

\noindent
The current \texttt{statevectorBackend} instantiation gives executable
semantics while preserving theorem-level invariants.  A key
normalization result ensures physical consistency:

\begin{theorem}[Post-Collapse Normalization]
\texttt{collapseByOutcome\_norm\_eq\_one\_of\_raw\_pos}:
if $0 < \texttt{rawOutcomeProb}(i, \mathit{out}, \mathit{st})$ then
$\|\psi'\| = 1$ after collapse.
\end{theorem}

\subsection{MQ-MORK Bridge}

MQ measurement branching is linked to MORK binary-fold branching
(\texttt{Interoperability.lean}, 37~lines) via:

\begin{theorem}[MQ-MORK Interoperability]
\texttt{comm\_nondeterminism\_iff\_mork\_binary}: the two COMM outcomes
of a measurement correspond bijectively to the two sub-results of
a MORK binary fold step.
\end{theorem}

\noindent
This ties communication-level nondeterminism in MQ directly into the
MORK three-phase execution model~\cite{MeredithStayOSLF2020}.


%% ====================================================================
\section{The MeTTa-Calculus}
\label{sec:metta-calc}
%% ====================================================================

The MeTTa-calculus~\cite{MeredithMeTTaCalculus2024} formalizes the core
of the MeTTa language as a symmetric reflective higher-order concurrent
calculus with backchaining.  The formalization comprises 9~files
totaling 904~lines under \texttt{ProcessCalculi/MeTTaCalculus/}.

\subsection{Syntax}

The process grammar is:
\[
  P,Q ::= \mathbf{0} \mid P \mathbin{|} Q \mid
           \mathsf{for}(t \leftarrow x)\, P \mid x\mathbin{?}P
           \mid {*}x
\]
\[
  x,y ::= @(P) \qquad
  t,u ::= (P) \mid \mathtt{true} \mid \mathtt{false} \mid n \mid s \mid \mathtt{sym}
\]
where $\mathsf{for}(t \leftarrow x)\, P$ is guarded input (listen for
pattern~$t$ on channel~$x$, then continue as~$P$), $x\mathbin{?}P$ is
reflection (inspect~$P$ and emit its future state on~$x$), $@(P)$
quotes a process into a name, and ${*}x$ drops (unquotes) a name back
to a process.

\texttt{Syntax.lean} (135~lines) encodes the grammar as abbreviations of
the shared \texttt{Pattern} AST:

\begin{lstlisting}[caption={MeTTa-calculus syntax (from \texttt{Syntax.lean})}]
abbrev Proc := Pattern   -- processes
abbrev Name := Pattern   -- names (quoted processes)
abbrev Term := Pattern   -- terms (data values)

def pZero    : Proc := .apply "MZero" []
def pPar     (elems : List Proc) : Proc :=
  .collection .hashBag elems none
def nQuote   (p : Proc) : Name := .apply "MQuote" [p]
def pDrop    (x : Name) : Proc := .apply "MDrop" [x]
def pFor     (t : Term) (x : Name) (k : Proc) : Proc :=
  .apply "MFor" [t, x, k]
def pReflect (x : Name) (p : Proc) : Proc :=
  .apply "MRef" [x, p]
\end{lstlisting}

\noindent
Two language definitions are provided: \texttt{mettaCalc} (full COMM +
REFL rules) and \texttt{mettaCalcCommOnly} (the COMM-only fragment, used
to avoid self-reference in the REFL rule's recursive stepping).

\subsection{Structural Congruence}

Structural congruence ($\equiv$) is the smallest equivalence closed
under:
\begin{align}
  P \mathbin{|} \mathbf{0}  &\equiv P
    \tag{\textsc{Nil}} \\
  P \mathbin{|} Q            &\equiv Q \mathbin{|} P
    \tag{\textsc{Comm}} \\
  (P \mathbin{|} Q) \mathbin{|} R
                              &\equiv P \mathbin{|} (Q \mathbin{|} R)
    \tag{\textsc{Assoc}} \\
  @({*}x)                    &\equiv x
    \tag{\textsc{QuoteDrop}}
\end{align}

\texttt{StructuralCongruence.lean} (67~lines) defines the inductive
relation \texttt{SC} (notation~$\approx_m$) with 11~constructors
covering all clauses above plus $\alpha$-equivalence, permutation, and
pointwise congruence:

\begin{lstlisting}[caption={Structural congruence (from \texttt{StructuralCongruence.lean})}]
inductive SC : Proc -> Proc -> Prop where
  | refl   (p : Proc) : SC p p
  | symm   (p q : Proc) : SC p q -> SC q p
  | trans  (p q r : Proc) : SC p q -> SC q r -> SC p r
  | alpha  (p q : Proc) : p = q -> SC p q
  | par_nil_left  (p : Proc) :
      SC (pPar [pZero, p]) p
  | par_nil_right (p : Proc) :
      SC (pPar [p, pZero]) p
  | par_comm  (p q : Proc) :
      SC (pPar [p, q]) (pPar [q, p])
  | par_assoc (p q r : Proc) :
      SC (pPar [pPar [p,q], r]) (pPar [p, pPar [q,r]])
  | par_perm  (xs ys : List Proc) :
      xs.Perm ys -> SC (pPar xs) (pPar ys)
  | par_cong  (xs ys : List Proc) :
      xs.length = ys.length ->
      (forall i hx hy, SC (xs.get {val:=i,isLt:=hx})
                           (ys.get {val:=i,isLt:=hy})) ->
      SC (pPar xs) (pPar ys)
\end{lstlisting}

\noindent
Paper-mapped equivalence theorems:
\begin{itemize}
\item \texttt{paper\_equiv\_par\_nil\_left}: $P \mathbin{|} \mathbf{0} \approx_m P$
\item \texttt{paper\_equiv\_par\_nil\_right}: $\mathbf{0} \mathbin{|} P \approx_m P$
\item \texttt{paper\_equiv\_par\_comm}: $P \mathbin{|} Q \approx_m Q \mathbin{|} P$
\item \texttt{paper\_equiv\_par\_assoc}: $(P \mathbin{|} Q) \mathbin{|} R \approx_m P \mathbin{|} (Q \mathbin{|} R)$
\end{itemize}

\subsection{Reduction Semantics}

The calculus has two reduction rules.  The \textsc{Comm} rule performs a
symmetric rendezvous: two guarded inputs listening on the same channel
unify their payload patterns and continue with substituted bodies:
\begin{equation}
  \frac{\sigma = \mathit{mgu}(t, u)}
       {\mathsf{for}(t \leftarrow x)\, P \mathbin{|}
        \mathsf{for}(u \leftarrow x)\, Q
        \;\longrightarrow\;
        P\dot\sigma \mathbin{|} Q\dot\sigma}
  \tag{\textsc{Comm}}
\end{equation}
where $\dot\sigma$ is \emph{dot-substitution}: each binding $(v \mapsto
T)$ in~$\sigma$ becomes $(v \mapsto @(T))$, quoting the term into a
name.  The \textsc{Refl} rule lets a process inspect its own future
via one-step \textsc{Comm}-only lookahead:
\begin{equation}
  \frac{P \longrightarrow_{\textsc{Comm}} P'}
       {x\mathbin{?}P \;\longrightarrow\;
        \mathsf{for}((P') \leftarrow x)\, \mathbf{0}}
  \tag{\textsc{Refl}}
\end{equation}

\texttt{Reduction.lean} (264~lines) encodes both rules as OSLF rewrite
rules with premise queries.  The \textsc{Comm} rule uses a
fuel-bounded first-order unifier (256~steps) with an occurs check:

\begin{lstlisting}[caption={COMM rewrite rule (from \texttt{Syntax.lean})}]
private def commSymRule : RewriteRule := {
  name := "CommSym"
  typeContext := [("t", .base "Term"), ("u", .base "Term"),
    ("x", .base "Name"), ("p", .base "Proc"),
    ("q", .base "Proc"), ("pOut", .base "Proc"),
    ("qOut", .base "Proc")]
  premises := [
    .relationQuery relMettaComm
      [.fvar "t", .fvar "u", .fvar "p", .fvar "q",
       .fvar "pOut", .fvar "qOut"] ]
  left := .collection .hashBag [
    .apply "MFor" [.fvar "t", .fvar "x", .fvar "p"],
    .apply "MFor" [.fvar "u", .fvar "x", .fvar "q"]
  ] (some "rest")
  right := .collection .hashBag
    [.fvar "pOut", .fvar "qOut"] (some "rest")
}
\end{lstlisting}

\begin{lstlisting}[caption={REFL rewrite rule (from \texttt{Syntax.lean})}]
private def reflRule : RewriteRule := {
  name := "Refl"
  typeContext := [("x", .base "Name"),
    ("p", .base "Proc"), ("pNext", .base "Proc")]
  premises := [.relationQuery relMettaStepNoReflect
    [.fvar "p", .fvar "pNext"]]
  left  := .apply "MRef" [.fvar "x", .fvar "p"]
  right := .apply "MFor"
    [.apply "MTermProc" [.fvar "pNext"],
     .fvar "x", .apply "MZero" []]
}
\end{lstlisting}

\noindent
The dot-substitution and unification implementations complete the
executable semantics:

\begin{lstlisting}[caption={Dot-substitution and step function (from \texttt{Reduction.lean})}]
def dotBindings (s : Bindings) : Bindings :=
  s.map (fun (x, v) => (x, nQuote v))

def applyDot (s : Bindings) (p : Proc) : Proc :=
  applyBindings (dotBindings s) p

def step (p : Proc) : List Proc :=
  rewriteWithContextWithPremisesUsing
    mettaCalcRelEnv mettaCalc p
\end{lstlisting}

\noindent
Positive and negative canaries validate both rules at the kernel level:
\begin{itemize}
\item \texttt{paper\_comm\_positive}/\texttt{paper\_comm\_negative}:
  a concrete COMM rendezvous succeeds; mismatched channels block.
\item \texttt{paper\_refl\_positive}/\texttt{paper\_refl\_negative}:
  a concrete REFL step emits a future state; a non-stepping body
  produces no output.
\end{itemize}

\subsection{Premise IR and Adequacy}

\texttt{Premises.lean} (88~lines) declares a \texttt{PremiseProgram}---a
datalog-style IR with two relations and two builtins:
\[
  \mathit{mettaComm}(t, u, P, Q, P', Q') \qquad
  \mathit{mettaStepNoReflect}(\mathit{src}, \mathit{tgt})
\]
\[
  \mathit{mettaCommWitness}(t, u, P, Q) \rightsquigarrow \mathit{MRef}(P', Q')
  \qquad
  \mathit{mettaCommOnlyStep}(\mathit{src}) \rightsquigarrow \mathit{tgt}
\]
with machine-checked stratification and wellformedness proofs.

\texttt{Adequacy.lean} (79~lines) bridges the abstract premise contract
to the concrete runtime:

\begin{theorem}[Premise--Runtime Adequacy]
\label{thm:metta-adequacy}
The following are proven:
\begin{enumerate}[nosep]
\item \texttt{runtime\_covers\_declared\_relations}: every declared
  relation is implemented by the runtime dispatcher.
\item \texttt{runtime\_covers\_declared\_builtins}: every declared
  builtin is implemented.
\item \texttt{stepViaPremiseContract\_eq\_step}: the IR-based
  stepping function agrees with the direct \texttt{step} function.
\item \texttt{reducesViaPremiseContract\_iff\_reduces}:
  $\texttt{ReducesViaPremiseContract}\;p\;q \iff \texttt{Reduces}\;p\;q$.
\end{enumerate}
\end{theorem}

\subsection{Interoperability with \texorpdfstring{$\rho$}{rho}-Calculus}

\texttt{Interoperability.lean} (113~lines) provides a shared-fragment
translation from the quote/drop/parallel subset of MeTTa-calculus into
$\rho$-calculus.  The translation \texttt{toRhoSharedProc?}~succeeds on
the translatable fragment and fails gracefully outside it.

For any successfully translated process, the file proves that the
$\rho$-image has a canonical open-map path-bisimulation witness:

\begin{itemize}
\item \texttt{shared\_translation\_pathBisim\_self}: the translated image
  is path-bisimilar to itself under the $\rho$-calculus
  \texttt{BisimulationKit}.
\item \texttt{shared\_translation\_esBisim\_self}: the corresponding
  $(E,S)$-bisimilarity corollary.
\item \texttt{paper\_shared\_to\_rho\_openmap\_self}: the paper-mapped
  wrapper.
\end{itemize}

\noindent
This establishes the shared-fragment translation as well-typed with
respect to the open-map framework, but is \emph{not yet} a full
cross-language bisimulation theorem.  The next step is to prove
forward/backward simulation for translated shared-fragment steps, then
upgrade to a full MeTTa$\leftrightarrow$$\rho$ bisimulation theorem.

\subsection{Regression and Paper Map}

\texttt{PaperMap.lean} (65~lines) provides a theorem-level index mapping
each clause of the specification~\cite{MeredithMeTTaCalculus2024} to its
Lean proof.
\texttt{Regression.lean} (65~lines) provides a regression corpus with
positive and negative test cases for both COMM and REFL rules.
\texttt{RelationNames.lean} (28~lines) is a single source of truth for
relation and builtin string identifiers.


%% ====================================================================
\section{Modal Mu-Calculus and the Rho-Mu Bridge}
\label{sec:mu-bridge}
%% ====================================================================

The modal $\mu$-calculus~\cite{Kozen1983,BradfieldStirling2007} extends
Hennessy-Milner logic with least and greatest fixed-point operators.
Our formalization covers standard semantics
(\texttt{ModalMuCalculus.lean}, 360~lines).

\subsection{Syntax}

Formulas use de~Bruijn indices for bound variables:
\[
  \phi ::= \top \mid \bot \mid \neg \phi \mid \phi \land \psi \mid
           \phi \lor \psi \mid \langle a \rangle \phi \mid [a] \phi
           \mid \mu X.\phi \mid \nu X.\phi \mid X
\]
where $\langle a \rangle$ is diamond (possibility), $[a]$ is box
(necessity), $\mu X$ is least fixed point, and $\nu X$ is greatest
fixed point.

\subsection{Rho-Calculus as LTS}

The $\rho$-calculus reduction relation $p \leadsto q$ is an LTS with
states = patterns, actions = $()$ (unlabeled), and transitions =
$\leadsto$.  This directly induces modal operators:

\begin{definition}[Reduction-Based Modalities]
\begin{align*}
\texttt{possiblyProp}\, \phi\, p &:= \exists q.\, p \leadsto q \land \phi\, q \quad \text{(diamond)} \\
\texttt{relyProp}\, \phi\, p &:= \forall q.\, q \leadsto p \to \phi\, q \quad \text{(box)}
\end{align*}
\end{definition}

\begin{theorem}[Galois Connection IS Modal Duality]
\label{thm:galois}
\[
(\forall p.\, \texttt{possiblyProp}\, \phi\, p \to \psi\, p) \iff
(\forall p.\, \phi\, p \to \texttt{relyProp}\, \psi\, p)
\]
This is exactly $\langle\rangle \phi \subseteq \psi \iff
\phi \subseteq [\,] \psi$.
\end{theorem}

\noindent
The proof is by direct unfolding of definitions
(\texttt{galois\_connection} in \texttt{Reduction.lean}).
This validates the OSLF construction~\cite{MeredithStayOSLF2020,leanOSLF2025}:
operational semantics naturally gives rise to modal logic.

\subsection{Spice Calculus as Iterated Diamond}

The Spice lookahead corresponds precisely to iterated modal operators:
\begin{align*}
  \texttt{presentMoment}(p) &\equiv \langle\rangle \phi \quad \text{(1-step)} \\
  \texttt{futureStates}(p, n) &\equiv \langle\rangle^n \phi \quad \text{($n$-step)} \\
  \texttt{reachableViaStarClosure}(p) &\equiv \nu X.\, \phi \lor \langle\rangle X \quad \text{(eventually)}
\end{align*}
All three equivalences hold by definitional equality (\texttt{simp}).
Precognitive agents computing Spice lookahead are computing
$\mu$-calculus temporal properties.

\subsection{Conjecture: Rho Strictly More Expressive Than Mu}
\label{sec:expressiveness}

\begin{conjecture}[$\rho$-Calculus Strictly More Expressive]
\label{conj:rho-mu-expressiveness}
There exists a property expressible in $\rho$-calculus that is not
expressible in modal $\mu$-calculus.  Specifically: $\mu$-calculus
cannot express ``this process has the capability to quote itself.''
\end{conjecture}

The intuition is that $\mu$-calculus formulas are determined entirely by
LTS behavior (which transitions occur), while $\rho$-calculus processes
can inspect their own syntactic structure via quote/unquote.  Two
processes $x[x]$ (self-quoting) and $y[z]$ (not self-quoting) are
behaviorally indistinguishable---both are dead---yet differ structurally.
No $\mu$-formula should be able to distinguish them.

A proper formalization would require: (1)~a coinductive bisimulation
relation (not just deadlock equivalence), (2)~a simulation theorem
showing $\rho$ can faithfully encode $\mu$-calculus evaluation, and
(3)~the separation argument over the actual $\rho$-calculus types from
\texttt{ProcessCalculi/RhoCalculus/}.  This remains future work.

\begin{remark}[G\"odel Analogy]
This is analogous to G\"odel's incompleteness theorem: first-order logic
cannot express ``this formula is provable''; adding a provability
predicate (reflection) gives strictly more expressive power.  Similarly,
$\mu$-calculus cannot express ``this process quotes itself'';
$\rho$-calculus can, via the quote operator $\mathsf{@}$.
\end{remark}


%% ====================================================================
\section{Quantale-Valued Semantics}
\label{sec:quantale}
%% ====================================================================

Standard $\mu$-calculus has Boolean semantics.  We generalize to
\emph{quantale-valued} semantics~\cite{Kozen1983} where satisfaction
takes values in a commutative quantale~$Q$
(\texttt{ModalQuantaleSemantics.lean}, 520~lines).

\begin{definition}[QLTS]
A \emph{quantale-labeled transition system} is a tuple
$(S, A, \mathrm{trans} : S \times A \times S \to Q)$ where $Q$ is a
commutative quantale with residuation.
\end{definition}

\begin{definition}[Quantale-Valued Satisfaction]
\begin{align*}
\texttt{qSatisfies}(s, \phi \land \psi) &= \texttt{qSatisfies}(s, \phi) \sqcap \texttt{qSatisfies}(s, \psi) \\
\texttt{qSatisfies}(s, \phi \lor \psi) &= \texttt{qSatisfies}(s, \phi) \sqcup \texttt{qSatisfies}(s, \psi) \\
\texttt{qSatisfies}(s, \langle a \rangle \phi) &= \bigsqcup_{s'} \mathrm{trans}(s, a, s') \otimes \texttt{qSatisfies}(s', \phi) \\
\texttt{qSatisfies}(s, [a] \phi) &= \bigsqcap_{s'} \mathrm{trans}(s, a, s') \Rightarrow \texttt{qSatisfies}(s', \phi) \\
\texttt{qSatisfies}(s, \mu X.\phi) &= \bigsqcap \{P \mid \texttt{transformer}(P) \sqsubseteq P\} \\
\texttt{qSatisfies}(s, \nu X.\phi) &= \bigsqcup \{P \mid P \sqsubseteq \texttt{transformer}(P)\}
\end{align*}
where $a \Rightarrow b := \bigsqcup\{z \mid z \otimes a \sqsubseteq b\}$
is left residuation.
\end{definition}

\begin{theorem}[Environment Monotonicity with Polarity]
\label{thm:env-mono}
For formula $\phi$ and variable $i$ with polarity
$p \in \{\mathrm{true}, \mathrm{false}\}$:
if $\rho_1(j) = \rho_2(j)$ for all $j \neq i$ and
$\rho_1(i) \sqsubseteq \rho_2(i)$, then
\[
  p = \mathrm{true} \implies \texttt{qSatisfies}(\rho_1, \phi)
    \sqsubseteq \texttt{qSatisfies}(\rho_2, \phi)
\]
\[
  p = \mathrm{false} \implies \texttt{qSatisfies}(\rho_2, \phi)
    \sqsubseteq \texttt{qSatisfies}(\rho_1, \phi)
\]
\end{theorem}

\noindent
The proof (130~lines) proceeds by structural induction on~$\phi$,
handling each constructor: base cases ($\top, \bot, X$) are immediate;
negation flips polarity via antitonicity of residuation;
conjunction/disjunction preserve polarity via $\sqcap$/$\sqcup$
monotonicity; diamond/box use multiplication and residuation
monotonicity; and the fixed-point cases ($\mu$/$\nu$) shift the de
Bruijn variable index, extend the environment, and show every
pre/post-fixed point for $\rho_2$ is also one for $\rho_1$.

This enables the Knaster-Tarski theorem, proving that fixed points
exist and are computable as limits of approximation sequences
(\texttt{muApprox\_mono}, \texttt{nuApprox\_antimono}).

\begin{example}[Instantiations]
\begin{itemize}[leftmargin=1.5em]
  \item \textbf{Probabilistic}: $Q = [0,1]$ with standard multiplication.
  \item \textbf{PLN evidence}: $Q = \mathbb{R}_{\geq 0}^\infty \times \mathbb{R}_{\geq 0}^\infty$ (positive/negative evidence counts).
  \item \textbf{Fuzzy}: $Q = [0,1]^n$ (multi-valued truth).
  \item \textbf{Tropical}: $Q = \mathbb{R}_{\geq 0}^\infty$ with $\min$/$+$ operations.
\end{itemize}
\end{example}


%% ====================================================================
\section{Common Infrastructure}
\label{sec:common}
%% ====================================================================

The three process calculi share approximately 40\% of their
reduction and congruence infrastructure.  A shared library in
\texttt{ProcessCalculi/Common/} (4~files, 246~lines) eliminates this
duplication through typeclasses and generic lemmas.

\subsection{Typeclasses}

\texttt{Common/ProcessAlgebra.lean} (31~lines) defines:

\begin{lstlisting}[caption={Shared process algebra typeclasses}]
class HasPar (P : Type) where
  par : P -> P -> P

class HasNu (P : Type) where
  nu : P -> P

class HasNil (P : Type) where
  nil : P
\end{lstlisting}

\noindent
\texttt{Common/Congruence.lean} (64~lines) adds \texttt{HasSC} for
structural congruence and the generic pattern
\texttt{reduces\_modulo\_sc}.

\subsection{Reflexive-Transitive Closure}

\texttt{Common/Star.lean} (148~lines) provides two flavors of
reflexive-transitive closure:
\begin{itemize}[leftmargin=1.5em]
  \item \texttt{RTClosure}: Type-valued (carries witness data).
  \item \texttt{RTClosureProp}: Prop-valued (for logical reasoning).
\end{itemize}
Generic lifters---\texttt{liftUnary}, \texttt{liftBinLeft},
\texttt{liftBinRight}---allow multi-step reduction to be lifted
through any unary or binary context constructor.

\subsection{Instance Registration}

Each calculus registers instances:
\begin{itemize}[leftmargin=1.5em]
  \item \textbf{$\pi$-calculus}: \texttt{HasPar}, \texttt{HasNil},
    \texttt{HasSC} (wrapping the Type-valued SC in \texttt{Nonempty}).
  \item \textbf{$\rho$-calculus}: \texttt{HasSC} (Prop-valued SC
    directly).
  \item \textbf{MQ-calculus}: \texttt{HasPar}, \texttt{HasNil},
    \texttt{HasNu}, \texttt{HasSC} (all four typeclasses).
\end{itemize}

\noindent
All three share: structural congruence as an equivalence relation with
parallel commutativity/associativity and scope extrusion (unified via
\texttt{HasSC}); and the COMM rule shape (sender + receiver $\to$
result), differing only in payload consumption (substitution in
$\pi$/$\rho$, measurement in MQ).


%% ====================================================================
\section{Summary and Future Work}
\label{sec:conclusion}
%% ====================================================================

\subsection{Summary}

We have presented a comprehensive Lean~4.28 formalization of five
process-algebraic systems and three cross-calculus bridges, totaling
17{,}551 lines with zero \texttt{sorry} and zero custom axioms
(Table~\ref{tab:summary}).

Key results include:
\begin{enumerate}[leftmargin=1.5em]
  \item The Galois connection in $\rho$-calculus \emph{is} modal
    diamond-box duality (Theorem~\ref{thm:galois}), validating the OSLF
    construction.
  \item Forward simulation of the $\pi$-to-$\rho$ encoding
    (Theorem~\ref{thm:forward-sim}) following
    Lybech~\cite{Lybech2022}, with 2{,}650~lines of backward admin
    reflection.
  \item Constructive both-outcomes witness for MQ measurement branching,
    with backend-parametric denotational semantics and post-collapse
    normalization.
  \item We conjecture that $\rho$-calculus is strictly more expressive
    than $\mu$-calculus (Conjecture~\ref{conj:rho-mu-expressiveness}):
    reflection cannot be captured by behavioral observation.
    A full formalization remains future work.
  \item Complete quantale-valued $\mu$-calculus semantics with
    environment monotonicity (Theorem~\ref{thm:env-mono}), enabling
    Knaster-Tarski fixed-point computation across probabilistic, fuzzy,
    tropical, and PLN-evidence domains.
  \item Spice calculus $n$-step lookahead is exactly iterated diamond
    $\langle\rangle^n$, proving that precognitive agents compute
    $\mu$-calculus temporal properties.
  \item MeTTa-calculus with executable COMM/REFL reduction, premise-IR
    adequacy bridge, and a shared-fragment translation into
    $\rho$-calculus with open-map self-bisimulation witness
    (Section~\ref{sec:metta-calc}).
\end{enumerate}

\subsection{Future Work}

\begin{enumerate}[leftmargin=1.5em]
  \item \textbf{Full backward simulation}: complete the
    $\pi$-to-$\rho$ backward simulation using the admin reflection
    infrastructure already formalized.
  \item \textbf{Multi-qubit MQ gates}: extend gate semantics beyond
    single-wire action; prove stronger collapse theory (preservation
    under composition, idempotence of repeated measurement).
  \item \textbf{PLN bridge completion}: formalize PLN's full inference
    calculus and prove temporal operators correspond to modal formulas
    with time-indexed LTS.
  \item \textbf{OSLF integration for MQ}: instantiate MQ processes as an
    OSLF language instance with $\Diamond \dashv \Box$ Galois
    connection, paralleling existing PeTTa and Governance instances.
  \item \textbf{Spice-MQ bridge}: connect $\rho$-calculus Spice
    $n$-step lookahead to MQ measurement prediction.
  \item \textbf{Hennessy-Milner theorem}: prove the classical result
    for the $\rho$-calculus LTS using \texttt{presentMoment\_finite}.
  \item \textbf{Full MeTTa$\leftrightarrow$$\rho$ bisimulation}: prove
    forward/backward simulation for translated shared-fragment steps,
    then upgrade to a full cross-language bisimulation theorem.
\end{enumerate}

\bibliographystyle{plain}
\bibliography{references}

\end{document}
